\documentclass{article}
\usepackage{mathtools} 
\usepackage{fontspec}
\usepackage[UTF8]{ctex}
\usepackage{amsthm}
\usepackage{mdframed}
\usepackage{xcolor}
\usepackage{amssymb}
\usepackage{amsmath}


% 定义新的带灰色背景的说明环境 zremark
\newmdtheoremenv[
  backgroundcolor=gray!10,
  % 边框与背景一致，边框线会消失
  linecolor=gray!10
]{zremark}{说明}

% 通用矩阵命令: \flexmatrix{矩阵名}{元素符号}{行数}{列数}
\newcommand{\flexmatrix}[4]{
  \[
  #1 = \begin{pmatrix}
    #2_{11}     & #2_{12}     & \cdots & #2_{1#4}   \\
    #2_{21}     & #2_{22}     & \cdots & #2_{2#4}   \\
    \vdots      & \vdots      & \ddots & \vdots     \\
    #2_{#31}    & #2_{#32}    & \cdots & #2_{#3#4}
  \end{pmatrix}
  \]
}

% 简化版命令（默认矩阵名为A，元素符号为a）: \quickmatrix{行数}{列数}
\newcommand{\quickmatrix}[2]{\flexmatrix{A}{a}{#1}{#2}}

\begin{document}
\title{习题 15.1}
\author{张志聪}
\maketitle

\section*{1}
\begin{itemize}
  \item (1)

        \begin{itemize}
          \item (i)

                $f(x)$在$[-\pi, \pi]$上是按段光滑的，可以展开为傅里叶级数
                \begin{align*}
                  a_0
                   & = \frac{1}{\pi} \int_{-\pi}^{\pi} f(x) dx \\
                   & = \frac{1}{\pi} \int_{-\pi}^{\pi} x dx    \\
                   & = \frac{1}{\pi} \int_{-\pi}^{\pi} x dx    \\
                   & = 0
                \end{align*}
                \begin{align*}
                  a_n
                   & = \frac{1}{\pi} \int_{-\pi}^{\pi} f(x) cos\, nx dx \\
                   & = \frac{1}{\pi} \int_{-\pi}^{\pi} x cos\, nx dx    \\
                   & = 0
                \end{align*}
                注：$x cos\, nx$是奇函数。
                \begin{align*}
                  b_n
                   & = \frac{1}{\pi} \int_{-\pi}^{\pi} f(x) sin\, nx dx                                   \\
                   & = \frac{1}{\pi} \int_{-\pi}^{\pi} x sin\, nx dx                                      \\
                   & = \frac{1}{\pi} (-\frac{1}{n})\int_{-\pi}^{\pi} x d(cos\, nx)                        \\
                   & = -\frac{1}{n\pi}(x cos\, nx \bigg|_{-\pi}^{\pi} - \int_{-\pi}^{\pi} cos\, nx \, dx) \\
                   & = -\frac{1}{n\pi}[\pi(-1)^n - (-\pi)(-1)^n - 0]                                      \\
                   & = (-1)^{n+1} \frac{2}{n}
                \end{align*}
                于是
                \begin{align*}
                  \frac{a_0}{2} + \sum\limits_{n = 1}^{\infty} a_n cos nx + b_n sin nx
                   & = \sum\limits_{n = 1}^{\infty} (-1)^{n+1} \frac{2}{n} sin nx
                \end{align*}
                由于$f(x)$在$(-\pi, \pi)$上连续，可知
                \begin{align*}
                  f(x) = \sum\limits_{n = 1}^{\infty} (-1)^{n+1} \frac{2}{n} sin nx , \, \, x \in (-\pi + 2k\pi, \pi + 2k\pi)
                \end{align*}
                由收敛定理可知，$x = \pm \pi$时
                \begin{align*}
                  \sum\limits_{n = 1}^{\infty} (-1)^{n+1} \frac{2}{n} sin nx
                   & = \frac{f(\pi + 0) + f(\pi - 0)}{2} \\
                   & = \frac{\pi + (-\pi)}{2}            \\
                   & = 0
                \end{align*}

          \item  (ii) 略
        \end{itemize}

  \item (2)

        \begin{itemize}
          \item (1)

                $f(x)$在$[-\pi, \pi]$上是按段光滑的，可以展开为傅里叶级数，有
                \begin{align*}
                  a_n
                   & = \frac{1}{\pi} \int_{-\pi}^{\pi} f(x) cos\, nx dx                                                     \\
                   & = \frac{1}{\pi} \int_{-\pi}^{\pi} x^2 cos\, nx dx                                                      \\
                   & = \frac{1}{n\pi} \int_{-\pi}^{\pi} x^2 d(sin\, nx)                                                     \\
                   & = \frac{1}{n\pi}\left( x^2 sin\, nx \bigg|_{-\pi}^{\pi} - \int_{-\pi}^{\pi} 2 x \, sin\, nx dx \right) \\
                   & = \frac{1}{n\pi}\left( 0 - \int_{-\pi}^{\pi} 2 x \, sin\, nx dx \right)                                \\
                   & = -\frac{2}{n\pi} \int_{-\pi}^{\pi} x \, sin\, nx dx                                                   \\
                   & = -\frac{2}{n\pi} (- \frac{1}{n}) \int_{-\pi}^{\pi} x \, d(cos\, nx)                                   \\
                   & = \frac{2}{n^2\pi}  \left( x \, cos\, nx \bigg|_{-\pi}^{\pi} - \int_{-\pi}^{\pi} cos\, nx dx \right)   \\
                   & = \frac{2}{n^2\pi}  \left[ (-1)^n 2\pi - 0 \right]                                                     \\
                   & = (-1)^n \frac{4}{n^2}
                \end{align*}
                \begin{align*}
                  a_0
                   & = \frac{1}{\pi} \int_{-\pi}^{\pi} f(x) dx                        \\
                   & = \frac{1}{\pi} \int_{-\pi}^{\pi} x^2 dx                         \\
                   & = \frac{1}{\pi} \left(\frac{1}{3}x^3 \bigg|_{-\pi}^{\pi} \right) \\
                   & = \frac{1}{\pi} \frac{2}{3} (\pi)^3                              \\
                   & = \frac{2}{3} (\pi)^2
                \end{align*}
                \begin{align*}
                  b_n
                   & = \frac{1}{\pi} \int_{-\pi}^{\pi} f(x) sin\, nx dx \\
                   & = \frac{1}{\pi} \int_{-\pi}^{\pi} x^2 sin\, nx dx  \\
                   & = 0
                \end{align*}
                故
                \begin{align*}
                  f(x) & =  \frac{a_0}{2} + \sum\limits_{n = 1}^{\infty} a_n cos nx + b_n sin nx                                \\
                       & = \frac{(\pi)^2}{3} \sum\limits_{n = 1}^{\infty} (-1)^n \frac{4}{n^2} cos \, nx , \, x \in (-\pi, \pi)
                \end{align*}
                \textcolor{red}{注：}
                $x = \pm \pi$时，
                \begin{align*}
                  \frac{f(\pi + 0) + f(\pi - 0)}{2}
                  = \pi^2
                   & = \frac{\pi^2}{3} + \sum\limits_{n = 1}^{\infty} (-1)^n \frac{4}{n^2} cos \, n\pi \\
                   & = \frac{\pi^2}{3} + \sum\limits_{n = 1}^{\infty} \frac{4}{n^2}
                \end{align*}
                经过变形，我们有
                \begin{align*}
                  \frac{\pi^2}{6} & = \sum\limits_{n = 1}^{\infty} \frac{1}{n^2}
                \end{align*}

          \item (2)

                $f(x)$在$[0, 2\pi]$上是按段光滑的，可以展开为傅里叶级数，有
                \begin{align*}
                  a_0 & = \frac{1}{\pi} \int_{0}^{2\pi} f(x) dx          \\
                  a_n & = \frac{1}{\pi} \int_{0}^{2\pi} f(x) cos\, nx dx \\
                  b_n & = \frac{1}{\pi} \int_{0}^{2\pi} f(x) sin\, nx dx
                \end{align*}
                计算略

        \end{itemize}

  \item (3)
        \begin{equation*}
          f'(x) = \begin{cases}
            a, & -\pi < x < 0 \\
            b, & 0 < x < \pi
          \end{cases}
        \end{equation*}
        又因为
        \begin{align*}
          \lim\limits_{x \to 0^-} f'(x) = a \\
          \lim\limits_{x \to 0^+} f'(x) = b
        \end{align*}
        所以$f(x)$在$[-\pi, \pi]$上是分段光滑的，可以展开为傅里叶级数
        \begin{align*}
          a_0
           & = \frac{1}{\pi} \int_{-\pi}^{\pi} f(x) dx                                        \\
           & = \frac{1}{\pi} \left(\int_{-\pi}^{0} ax \, dx + \int_{0}^{\pi} bx \, dx \right) \\
           & = \frac{b - a}{2} \pi
        \end{align*}
        \begin{align*}
          a_n
           & = \frac{1}{\pi} \int_{-\pi}^{\pi} f(x) cos\, nx dx                                                 \\
           & = \frac{1}{\pi} \left(\int_{-\pi}^{0} ax cos\, nx \, dx + \int_{0}^{\pi} bx cos\, nx \, dx \right) \\
           & = \frac{a - b}{n^2 \pi} [1 - (-1)^n]
        \end{align*}
        \begin{align*}
          b_n
           & = \frac{1}{\pi} \int_{-\pi}^{\pi} f(x) sin\, nx dx                                                 \\
           & = \frac{1}{\pi} \left(\int_{-\pi}^{0} ax sin\, nx \, dx + \int_{0}^{\pi} bx sin\, nx \, dx \right) \\
           & = (-1)^n \frac{a + b}{n}
        \end{align*}
        故
        \begin{align*}
           & \frac{a_0}{2} + \sum\limits_{n = 1}^{\infty} a_n cos nx + b_n sin nx \\
           & = \frac{b - a}{4} \pi +
          \sum\limits_{n = 1}^{\infty} \frac{a - b}{n^2 \pi} [1 - (-1)^n] cos nx +  (-1)^n \frac{a + b}{n} sin nx
        \end{align*}
\end{itemize}

\section*{2}

提示：利用第九章$\S 5$习题6；

% \section*{3}

% 显然$f$是分段光滑的，可以展开为傅里叶级数
% \begin{align*}
%   a_0 & = \frac{1}{\pi} \int_{-\pi}^{\pi} f(x) dx                                                          \\
%       & = \frac{1}{\pi} \int_{-\pi}^{0} - \frac{\pi}{4} dx + \frac{1}{\pi} \int_{0}^{\pi} \frac{\pi}{4} dx \\
%       & = 0
% \end{align*}

\end{document}